% Supplementary Material for PAPER1_FREEFREE_DRAFT.tex
% Mapping (main text uses static ``Supplementary Material, §S.n'' pointers):
%   S.1  former App. A (sec:frobenius-appendix)
%   S.2  former main-text appendix (B.1 + Green) + (n,s) algebra
%        (B.3 / B.3$' already lived here; no remaining main-text appendix)
%   S.3  former App. C minus tab:shi (promoted into main-text §6.2)
%   S.3.1 / S.3.2 / S.3.3 = conv / instruments / cutoff
%   S.3.4 / S.3.5 = extensional close-pairs / Shi OOP isolation
%   S.3.6 = raw-to-physical conversion constants (moved from main-text §6.12)
%   S.3.7 = SUBDOM numerical procedure and gate detail
%   Tables S.1--S.8 sit in S.3 (S.8 = flexural spurious zeros, was main Table 7)
%   S.4  cluster job-number provenance (added 2026-08-25, advisor pass)
\documentclass[11pt]{article}
\usepackage[margin=1in,letterpaper]{geometry}
\usepackage{amsmath,amssymb,booktabs,array,xcolor}
\usepackage[hidelinks]{hyperref}
\usepackage{placeins}
% Match the main draft's float placement so tall tables (e.g. Table S.2) sit at
% the top of a text page rather than claiming a float page of their own.
\renewcommand{\topfraction}{0.9}
\renewcommand{\bottomfraction}{0.85}
\renewcommand{\textfraction}{0.07}
\renewcommand{\floatpagefraction}{0.85}
\setcounter{topnumber}{4}
\setcounter{bottomnumber}{4}
\setcounter{totalnumber}{8}
\emergencystretch=2.5em

\renewcommand{\thesection}{S.\arabic{section}}
\renewcommand{\thesubsection}{S.\arabic{section}.\arabic{subsection}}
\renewcommand{\thetable}{S.\arabic{table}}
\renewcommand{\theequation}{S.\arabic{equation}}

\title{\bfseries Supplementary Material:\\
Exact Wave-Function Solution and Spurious-Mode Discrimination
for Completely Free Annular Sector Plates}
\author{Garrek McCune$^{a,*}$, Daniel Stutts$^{a}$\\
\small\itshape $^a$Mechanical and Aerospace Engineering,\\
Missouri University of Science and Technology, Rolla, MO 65409, USA\\
\small\upshape $^*$Corresponding author. E-mail: ghmkfh@umsystem.edu}
\date{DRAFT --- \today}

\begin{document}
\maketitle

\noindent
This file is the complete Supplementary Material to the companion
manuscript: §§S.1--S.4, Tables S.1--S.8, and the three references
cited below. Equation tags (B.1), (B.2), (B.3) and (B.3$'$) retain the
source-paper numbering. Main-text \S3.3 cites (B.1), (B.2) and
(B.3$'$); (B.3) appears only here.

\section{Frobenius series construction}\label{sec:frobenius-appendix}
Following Ref.~\cite{seok1} \S3 (out-of-plane) and Ref.~\cite{seok2} \S3
(in-plane), the radial ordinary differential equation is expanded in
Frobenius series about the sector's center radius $r_0$ (equivalently
$r^{*}=0$) rather than about the singular point $r=0$, so that all four
independent solutions remain regular on $r^{*}\in[-\pi/2,\pi/2]$:
\[
G_q(r^{*},\xi)=\sum_{m=0}^{M}a^{*(q)}_m(\xi)\,r^{*\,\lambda(q)+m},
\qquad q=1,\dots,4 \quad\text{(Ref.~\cite{seok1} Eq.~(32))}.
\]
The indicial equation,
\[
\bar r_0^{4}\,a^{*}_0\,\lambda(\lambda-1)(\lambda-2)(\lambda-3)=0
\quad\text{(Ref.~\cite{seok1} Eq.~(31))},
\]
fixes $\lambda(q)=q-1$ for $q=1,\dots,4$, giving the four out-of-plane
branches their characteristic near-center behavior ($r^{*0},r^{*1},r^{*2},
r^{*3}$). The paper states the resulting recursion for the coefficients
$a^{*(q)}_m$ explicitly only through this indicial step, noting it is
``much too cumbersome to include'' in closed form (Ref.~\cite{seok1} \S3);
the present implementation carries the recursion numerically to
convergence at working precision (\texttt{mpmath}, 40 digits) for each
candidate frequency, rather than truncating at a fixed low order the way
the papers' fixed-precision (20-significant-digit Maple) calculation could
tolerate. Truncation order $M$ is chosen so that successive terms no longer
change the retained digits of the radial function at the working
precision; no fixed $M$ is hardcoded independent of $\text{dps}$.

The in-plane problem is analogous but vector-valued: the coupled radial
functions $F,G$ expand as
\[
F(r^{*},\zeta)=\sum_{m=0}^{M}a^{*(q)}_m(\zeta)\,r^{*\,\lambda(q)+m},\qquad
G(r^{*},\zeta)=\sum_{m=0}^{M}b^{*(q)}_m(\zeta)\,r^{*\,\lambda(q)+m}
\quad\text{(Ref.~\cite{seok2} Eqs.~(28)--(29))},
\]
with indicial equations
\[
a^{*}_0\,\bar{c}^{*}_{11}\,\bar r_0\,\lambda(\lambda-1)=0,\qquad
b^{*}_0\,\bar r_0\,\lambda(\lambda-1)=0 \quad\text{(Ref.~\cite{seok2}
Eq.~(30a,b); \(\bar{c}^{*}_{11}=c^{*}_{11}/c_{66}\))},
\]
giving $\lambda=0,1$ (each a double root of the coupled system) and four
coupled solutions selected by the uncoupled leading-coefficient scheme of
Ref.~\cite{seok2} Eq.~(32a--d): the two solutions built on $\lambda=0$
differ in which of $F,G$ carries the non-vanishing leading coefficient, and
likewise for $\lambda=1$, giving four independent, non-degenerate coupled
radial functions overall.

In both the out-of-plane and in-plane cases, the transverse wavenumber
($\xi$ or $\zeta$) that indexes a physically admissible solution is a root
of the radial-edge dispersion relation, and as the frequency $\Omega$
varies this root traces a curve that passes through real, imaginary, and
complex regimes (Ref.~\cite{seok1} Fig.~2 plots several such branches).
The branch-census construction tracks every branch continuously across
frequency and re-fills any family a purely local root-finder would drop
at a branch collision, so that the retained basis at every $\Omega$
contains one representative from each of the four families with no
silent gaps, the practically difficult part of the replication,
independent of the recursion detail above.

\section{Free-edge surface-term algebra}\label{sec:surface-algebra}
\label{sec:corner-appendix}
\paragraph{The cantilever's two corner terms (Eq.~(43)).} The tail of
Ref.~\cite{seok1} Eq.~(43), confirmed directly against the published
page rather than inferred, reads
\begin{equation}
+(T-\bar\nu)\left[\frac{\partial^2}{\partial r^{*}\partial\theta}
\left\{\frac{w}{r^{*}+\bar r_0}\right\}\right]_{r^{*}=-\pi/2,\,\theta=-\Theta}^{r^{*}=\pi/2,\,\theta=-\Theta}
\;-\;2(T-\bar\nu)\left[\frac{\partial^2}{\partial r^{*}\partial\theta}
\left\{\frac{w}{r^{*}+\bar r_0}\right\}\right]_{r^{*}=-\pi/2,\,\theta=\Theta}^{r^{*}=\pi/2,\,\theta=\Theta}
=0, \tag{B.1}
\end{equation}
appended to an integral term (over $r^{*}\in[-\pi/2,\pi/2]$, evaluated at
$\theta=-\Theta$) that enforces the clamped edge's kinematic conditions in
mixed/Lagrangian form, and a boundary term at $\theta=\Theta$ that gives
the free edge's own moment condition. Both displayed corner terms are jump
contributions over the two radial edges $r^{*}=\pm\pi/2$ (i.e., $r=R_o$
minus $r=R_i$), but at a single, different circumferential edge apiece: the
$+(T-\bar\nu)$ term at the \emph{clamped} edge $\theta=-\Theta$, and the
$-2(T-\bar\nu)$ term at the \emph{free} edge $\theta=\Theta$. The two
coefficients are not a sign-flipped pair of one functional: the clamped-edge
term is generated by the paper's substitution enforcing $w=\partial
w/\partial\theta=0$ along $\theta=-\Theta$ (a constraint on the primal
solution, entering the variational statement through the same
Lagrangian-substitution mechanism as the integral term above it), while the
free-edge term is a genuine natural boundary contribution from the free
edge's own variation. Consequently the free-free extension's second free
edge ($\theta=-\Theta$, now released rather than clamped) must reuse the
\emph{free}-edge formula (the $-2(T-\bar\nu)$ structure, correctly
re-oriented) and cannot reuse Eq.~(43)'s $+(T-\bar\nu)$ coefficient,
which is specific to a kinematically clamped edge and has no meaning once
that edge's kinematic constraint is removed.

The asymmetry is not an arbitrary bookkeeping choice: it has a classical-plate-theory origin. The free edge's factor of 2 is the same
doubling that produces the familiar concentrated corner reaction
$2M_{r\theta}$ at a free corner in Kirchhoff plate theory, where the two
adjoining free edges' Kirchhoff shear conditions combine into a single
term; the clamped edge has no analogous free-variation contribution to
double, since its kinematic condition removes that degree of freedom from
the variation entirely, leaving the single-multiplier Lagrangian term
instead. The link from
Ref.~\cite{seok1}'s already-published general functional (Eq.~(6)) to this
specific free-free result is derived explicitly, not asserted by analogy,
in the algebra below; main-text \S3.3 states the resulting free-free
corner term and proves the same-parity cancellation of the mis-oriented
assembly. What is genuinely not
reproduced here, and is a different and larger undertaking than the above,
is an independent line-by-line integration-by-parts derivation of
Eq.~(6)/Eq.~(1a) itself starting from the governing equation: that is
Ref.~\cite{seok1}'s own foundational derivation, already published and
peer-reviewed, and identical for the cantilever and free-free cases alike,
so redoing it here would not add evidence specific to the free-free
extension.

The intermediate algebra that produces (B.3) and (B.3$'$) from
Ref.~\cite{seok1} Eq.~(6) follows.

\paragraph{Grounding the corner-term asymmetry in the general functional (Eq.~(6)).}
Ref.~\cite{seok1} Eq.~(6), confirmed directly against the published
page rather than inferred, is the true pre-specialization variational
statement from which Eqs.~(42)--(43) are later obtained by substituting
the cantilever's specific edge conditions:
\begin{align}
&\int_S\!\Big[\Big\{\tfrac{\partial\tau^{(0)}_{rz}}{\partial r}
+\tfrac1r\Big(\tfrac{\partial\tau^{(0)}_{\theta z}}{\partial\theta}
+\tau^{(0)}_{rz}\Big)-2\rho h\ddot w\Big\}\delta w
+\Big\{\tfrac{\partial\tau^{(1)}_{rr}}{\partial r}
+\tfrac1r\tfrac{\partial\tau^{(1)}_{r\theta}}{\partial\theta}
+\tfrac{\tau^{(1)}_{rr}-\tau^{(1)}_{\theta\theta}}{r}
-\tau^{(0)}_{rz}\Big\}\delta u_r^{(1)} \notag\\
&\quad+\Big\{\tfrac{\partial\tau^{(1)}_{r\theta}}{\partial r}
+\tfrac1r\Big(\tfrac{\partial\tau^{(1)}_{\theta\theta}}{\partial\theta}
+2\tau^{(1)}_{r\theta}\Big)-\tau^{(0)}_{\theta z}\Big\}\delta u_\theta^{(1)}\Big]\,dS
-\int_{-\Theta}^{\Theta}\!\Big[r\Big\{\big(\tau^{(0)}_{rz}
+\tfrac1r\tfrac{\partial\tau^{(1)}_{r\theta}}{\partial\theta}\big)\delta w
-\tau^{(1)}_{rr}\,\delta\big(\tfrac{\partial w}{\partial r}\big)\Big\}\Big]_{r=R_i}^{r=R_o}d\theta \notag\\
&\quad-\int_{R_i}^{R_o}\!\Big[w\,\delta\tau^{(0)}_{\theta z}
-\tfrac{\partial w}{\partial r}\delta\tau^{(1)}_{r\theta}\Big]_{\theta=-\Theta}dr
-\int_{R_i}^{R_o}\!\Big[\big(\tau^{(0)}_{\theta z}
+\tfrac{\partial\tau^{(1)}_{\theta r}}{\partial r}\big)\delta w
-\tau^{(1)}_{\theta\theta}\,\delta\big(\tfrac1r\tfrac{\partial w}{\partial\theta}\big)\Big]_{\theta=\Theta}dr \notag\\
&\quad-\big[\tau^{(1)}_{r\theta}\delta w\big]^{r=R_o,\,\theta=-\Theta}_{r=R_i,\,\theta=-\Theta}
+2\big[\tau^{(1)}_{r\theta}\delta w\big]^{r=R_o,\,\theta=\Theta}_{r=R_i,\,\theta=\Theta}=0.
\tag{B.2, $\equiv$ Ref.~\cite{seok1} Eq.~(6)}
\end{align}
Immediately after Eq.~(6), Ref.~\cite{seok1} states that its last two
surface terms were obtained ``by performing the line integral around the
contour $C_N$ in Eq.~(1a) along the three free edges'' -- i.e., Eq.~(6) as
printed is \emph{already} built for a specific edge assignment (both radial
edges plus $\theta=\Theta$ free; $\theta=-\Theta$ not among the ``three
free edges''), not a symmetric statement later specialized equally at both
circumferential edges. This is visible directly in the two circumferential
boundary terms' structure: the $\theta=\Theta$ term is a stress (dual)
quantity multiplying a variation of the primal deflection ($\delta w$,
$\delta(\tfrac1r\partial w/\partial\theta)$) -- the natural form whose
stationarity under arbitrary $\delta w$ gives the free edge's own moment
and shear conditions -- while the $\theta=-\Theta$ term instead multiplies
the primal deflection ($w$, $\partial w/\partial r$) against a variation of
the dual/stress quantities ($\delta\tau^{(0)}_{\theta z}$,
$\delta\tau^{(1)}_{r\theta}$), the mixed/Lagrangian form whose
stationarity under arbitrary $\delta\tau$ instead constrains the primal
kinematic quantities themselves. This is the level at which Eq.~(B.1)'s
coefficient asymmetry actually originates: not in Eq.~(43)'s corner
coefficients as such, but in Eq.~(6) treating $\theta=-\Theta$ and
$\theta=\Theta$ asymmetrically from the outset, before any constitutive
substitution.

For the free-free extension, releasing $\theta=-\Theta$ from its clamped
condition means its Eq.~(6) boundary term must be replaced by the natural
form. This substitution is carried out here explicitly from Eq.~(1a)'s
general indicial contour term $n_a\tau^{(1)}_{a\beta}s_\beta$, rather than
asserted by pattern-matching, by transforming it into $r,\theta$ components
at $\theta=\Theta$ and $\theta=-\Theta$ and tracking the outward-normal/
tangent reversal between them through a counterclockwise traversal of the
sector boundary starting at $(R_i,-\Theta)$: the $\theta=-\Theta$ edge is
traversed with increasing $r$ (tangent $\hat s=+\hat r$, outward normal
$\hat n=-\hat\theta$), the $\theta=\Theta$ edge with decreasing $r$ (tangent
$\hat s=-\hat r$, outward normal $\hat n=+\hat\theta$). At $\theta=\Theta$
($n=+\hat\theta$, $s=-\hat r$): $n_a t^{(0)}_{a3}=\tau^{(0)}_{\theta z}$,
$u^{(0)}_{3,n}=\tfrac1r\partial w/\partial\theta$,
$n_a t^{(1)}_{a\beta}s_\beta=-\tau^{(1)}_{\theta r}$, so
$(n_a t^{(1)}_{a\beta}s_\beta)_{,s}=\partial\tau^{(1)}_{\theta r}/\partial r$
(the $s$-derivative along $-\hat r$ flips the sign of the $r$-derivative),
and $n_a t^{(1)}_{a\beta}n_\beta=\tau^{(1)}_{\theta\theta}$; substituting
into Eq.~(1a)'s third term and converting $ds\to dr$ (tangent direction and
integration limits both reverse, canceling) reproduces Eq.~(6)'s printed
$\theta=\Theta$ bracket exactly, leading minus sign included -- a working
self-consistency check before the method is trusted at the untested edge.
At $\theta=-\Theta$ ($n=-\hat\theta$, $s=+\hat r$):
$n_a t^{(0)}_{a3}=-\tau^{(0)}_{\theta z}$,
$u^{(0)}_{3,n}=-\tfrac1r\partial w/\partial\theta$,
$n_a t^{(1)}_{a\beta}s_\beta=-\tau^{(1)}_{\theta r}$ (the same value as at
$\Theta$, since for a symmetric tensor $n_a\tau_{a\beta}s_\beta$ is
invariant under simultaneously flipping both $n\to-n$ and $s\to-s$ -- a
useful independent check), so
$(n_a t^{(1)}_{a\beta}s_\beta)_{,s}=-\partial\tau^{(1)}_{\theta r}/\partial r$
(no sign flip this time, since $s=+\hat r$ already matches increasing $r$),
and $n_a t^{(1)}_{a\beta}n_\beta=\tau^{(1)}_{\theta\theta}$ (two minus signs
cancel). The integrand becomes the \emph{negative} of the $\theta=\Theta$
bracket, and with $ds=+dr$ (no limit-flip this time) Eq.~(1a)'s leading
minus combines with this internal minus to give an overall plus: the
\emph{same natural-edge bracket} as the $\theta=\Theta$ term, with an
\emph{opposite overall sign},
\begin{equation}
+\int_{R_i}^{R_o}\!\Big[\big(\tau^{(0)}_{\theta z}
+\tfrac{\partial\tau^{(1)}_{\theta r}}{\partial r}\big)\delta w
-\tau^{(1)}_{\theta\theta}\,\delta\big(\tfrac1r\tfrac{\partial w}{\partial\theta}\big)\Big]_{\theta=-\Theta}dr,
\tag{B.3}
\end{equation}
i.e. Eq.~(6)'s $\theta=\pm\Theta$ terms become
$\big[(\cdot)\delta w-(\cdot)\delta(\cdot)\big]_{-\Theta}^{\Theta}$ in the
Eq.~(7) bracket sense. This sign-alternation is confirmed two independent
ways: by the ($n,s$)-orientation bookkeeping above, and by the assembled
free-edge integral, which uses the same stress/displacement bracket on
every edge with orientation $+1$ at $\theta=\Theta$ and $-1$ at
$\theta=-\Theta$: exactly ``same bracket, opposite sign.''

Fully substituting Eq.~(B.3) into explicit deflection terms uses the same
two constitutive relations Ref.~\cite{seok1} applies at $\theta=\Theta$
(Eq.~(9c) and Eq.~(15) enter only the clamped edge's superseded
\emph{mixed}-term substitution, not this one), transcribed directly from
the source (Ref.~\cite{seok1} pp.~763--764):
\begin{align}
\tau^{(1)}_{\theta\theta}&=-\bar D\Big\{\bar\nu\frac{\partial^2w}{\partial r^2}
+\frac{R^{*}}{r}\Big(\frac{\partial w}{\partial r}+\frac1r\frac{\partial^2w}{\partial\theta^2}\Big)\Big\},
\tag{Ref.~\cite{seok1} Eq.~(14)}\\
\tau^{(0)}_{\theta z}+\frac{\partial\tau^{(1)}_{\theta r}}{\partial r}
&=-\bar D\Big\{\frac{R^{*}}{r^3}\frac{\partial^3w}{\partial\theta^3}
-\frac{(2T-R^{*})}{r^2}\frac{\partial^2w}{\partial r\partial\theta}
+\frac{2T}{r}\Big(\frac{\partial^3w}{\partial r^2\partial\theta}
+\frac1{r^2}\frac{\partial w}{\partial\theta}\Big)\notag\\
&\qquad-\bar\nu\Big(\frac1r\frac{\partial^3w}{\partial r^2\partial\theta}
+\frac2{r^3}\frac{\partial w}{\partial\theta}
-\frac2{r^2}\frac{\partial^2w}{\partial r\partial\theta}\Big)\Big\},
\tag{Ref.~\cite{seok1} Eq.~(20)}
\end{align}
together with the nondimensionalization $r=(2b/\pi)(r^{*}+\bar r_0)$
(Ref.~\cite{seok1} Eqs.~(26)--(28)), so $\partial/\partial
r=(\pi/2b)\,\partial/\partial r^{*}$ and $1/r=(\pi/2b)/(r^{*}+\bar r_0)$
($T,R^{*},\bar\nu$ are pure material-constant ratios, unaffected by this
substitution; $\theta$ itself is not rescaled). Substituting Eqs.~(14) and
(20) into the $\theta=\Theta$ bracket and factoring reproduces
Ref.~\cite{seok1}'s own published $\theta=\Theta$ block of Eq.~(43)
term-for-term (a self-consistency check with no reference to that
printed result used in the substitution itself), and, since Eqs.~(14) and
(20) hold pointwise at any $(r,\theta)$ with no reference to $\Theta$ or
$-\Theta$, the identical algebra applies verbatim at $\theta=-\Theta$ with
only Eq.~(B.3)'s outer sign flipped. The complete, explicit result, combining both edges in the
Eq.~(7) bracket sense, is
\begin{align}
&\int_{-\pi/2}^{\pi/2}\!dr^{*}\bigg[\frac1{r^{*}+\bar r_0}\frac{\partial}{\partial\theta}
\Big\{\frac{2Tw+R^{*}\partial_\theta^2w}{(r^{*}+\bar r_0)^2}
+2T\partial_{r^{*}}^2w-\frac{(2T-R^{*})}{r^{*}+\bar r_0}\partial_{r^{*}}w\notag\\
&\qquad-\bar\nu\Big\{\frac{2w}{(r^{*}+\bar r_0)^2}
+(r^{*}+\bar r_0)^2\partial_{r^{*}}\Big(\frac{\partial_{r^{*}}w}{(r^{*}+\bar r_0)^2}\Big)\Big\}\Big\}\,\delta w\notag\\
&\qquad-\Big\{\frac{R^{*}}{r^{*}+\bar r_0}\Big(\frac{\partial_\theta^2w}{(r^{*}+\bar r_0)^2}
+\frac{\partial_{r^{*}}w}{r^{*}+\bar r_0}\Big)
+\frac{\bar\nu}{r^{*}+\bar r_0}\partial_{r^{*}}^2w\Big\}\,\delta(\partial_\theta w)\bigg]_{\theta=-\Theta}^{\theta=\Theta}=0,
\tag{B.3$'$}
\end{align}
every term reached by explicit substitution of Eqs.~(14), (20) and the
nondimensionalization above at both edges, not reused at $\theta=-\Theta$
by sign-flip analogy from the published $\theta=\Theta$ block.

\paragraph{Closed-contour corner jump.} For the free-free case the
contour $C_N$ is the full closed boundary, and the corner coefficient
follows from Eq.~(1a)'s fourth term,
$\int_{c_N}\!ds\,(n_a\tau^{(1)}_{a\beta}s_\beta\,\delta w)_{,s}$, an
exact differential along arc length whose segment-by-segment integration
over a closed contour does not vanish, because $n_a,s_b$ are
discontinuous at each corner even though the physical stress field
$\tau^{(1)}_{r\theta}$ is continuous there. Writing
$F\equiv n_a\tau^{(1)}_{a\beta}s_\beta$ (using each segment's own
$(n,s)$ convention) and summing the fundamental-theorem contributions
corner by corner for the four segments traversed counterclockwise from
$(R_i,-\Theta)$:
\begin{center}
\begin{tabular}{lccc}
\toprule
segment & tangent $\hat s$ & normal $\hat n$ & $F=n_a\tau^{(1)}_{a\beta}s_\beta$\\
\midrule
$r=R_o$ arc ($\theta:-\Theta\to\Theta$) & $+\hat\theta$ & $+\hat r$ & $+\tau^{(1)}_{r\theta}$\\
$\theta=\Theta$ radial ($r:R_o\to R_i$) & $-\hat r$ & $+\hat\theta$ & $-\tau^{(1)}_{r\theta}$\\
$r=R_i$ arc ($\theta:\Theta\to-\Theta$) & $-\hat\theta$ & $-\hat r$ & $+\tau^{(1)}_{r\theta}$\\
$\theta=-\Theta$ radial ($r:R_i\to R_o$) & $+\hat r$ & $-\hat\theta$ & $-\tau^{(1)}_{r\theta}$\\
\bottomrule
\end{tabular}
\end{center}
and taking the jump at each corner as (incoming segment's $F$) minus
(outgoing segment's $F$),
\begin{align*}
(R_o,-\Theta):&\ (-\tau_{r\theta})-(\tau_{r\theta})=-2\tau_{r\theta}\big|_{(R_o,-\Theta)}, &
(R_o,\Theta):&\ (\tau_{r\theta})-(-\tau_{r\theta})=+2\tau_{r\theta}\big|_{(R_o,\Theta)},\\
(R_i,\Theta):&\ (-\tau_{r\theta})-(\tau_{r\theta})=-2\tau_{r\theta}\big|_{(R_i,\Theta)}, &
(R_i,-\Theta):&\ (\tau_{r\theta})-(-\tau_{r\theta})=+2\tau_{r\theta}\big|_{(R_i,-\Theta)}.
\end{align*}
Summing the $\theta=\Theta$ pair reproduces Ref.~\cite{seok1}'s own printed
cantilever corner term exactly,
$(R_o,\Theta)+(R_i,\Theta)=2\big[\tau^{(1)}_{r\theta}\delta w\big]^{R_o,\Theta}_{R_i,\Theta}$,
a self-consistency check with no reference to the deployed code; and
summing the $\theta=-\Theta$ pair, now smooth-turn corners of the closed
contour rather than the cantilever's open ends, gives
$(R_o,-\Theta)+(R_i,-\Theta)=-2\big[\tau^{(1)}_{r\theta}\delta w\big]^{R_o,-\Theta}_{R_i,-\Theta}$:
magnitude 2 at every corner, alternating in sign between the two edge
pairs. Substituting
$\tau^{(1)}_{r\theta}=-\bar D(T-\bar\nu)\partial^2(w/r)/\partial r\partial\theta$
(Ref.~\cite{seok1} Eq.~(15)) and the same nondimensionalization used for
Eq.~(B.3$'$) produces the free-free analogue of Eq.~(43) stated in
main-text \S3.3.

\paragraph{Independent confirmation via Green's identity.} The
closed-contour jump just given fixes the corner term's magnitude (2, every
corner) and its alternation pattern between the two edge pairs, but not,
by itself, which edge carries the ``$+$'' label in an isolated equation,
an overall sign-convention choice rather than physical content. A third,
code- and FE-independent route settles it: a Green's-identity argument on
the internal virtual-work bilinear form, evaluated on the classical
constant pure-twist field ($M_{r\theta}=\text{const}$,
$M_{rr}=M_{\theta\theta}=0$, which satisfies the free-edge natural
conditions identically), fixes the corner force to $R=2M_{ns}$ with its
sign read directly from the outgoing edge under $\partial\Omega$'s
standard positive (counterclockwise) orientation; applied to the $(n,s)$
table above, it reproduces the free-edge magnitude and the
factor of 2. Magnitude and alternation therefore follow from the contour
and Green arguments. Which isolated edge carries the ``$+$'' label is a
convention of the overall stationarity equation --- multiplying it by $-1$
swaps both labels and moves no frequency --- and what is required instead,
namely that the corner term and the surface term carry the same $\pm1$
orientation factor within one equation, is what the deployed assembly
implements (main-text \S3.3).

The free-free corner term is therefore constrained three ways (the
closed-contour jump of this section, the Green's-identity argument
just above, and comparison against the deployed assembly) rather than
asserted by pattern-matching against Eq.~(43) alone. The term-by-term
factoring behind Eq.~(B.3$'$) is the algebra of this section. The
same-parity cancellation of the mis-oriented assembly is proved in
main-text \S3.3.

\section{Additional validation tables}\label{sec:extra-tables}
The tables and supporting accounts below are referenced from
main-text \S4--\S6 but collected here so the main validation narrative
can stay on the tables that carry the paper's primary claims.
§S.3.6 is the raw-to-physical conversion algebra; §S.3.7 is the SUBDOM
numerical procedure; Table~S.8 is the
flexural spurious-zero listing moved from the main text.

\subsection{Basis-size convergence, per mode}\label{sec:conv-appendix}
The main text's \S4 summary states that 24 of 34 tabulated modes change by
less than $0.1\%$ across $n=20$, $28$, $36$, with the remainder traced to
two understood mechanisms. Full detail: nine labels (twelve (label, $n$)
cells, three labels flagged at two basis sizes each) initially returned a
value 0.6--12\% away from their own seed at one basis size; re-polishing
exactly those cells inside a bracket chosen to exclude the specific
neighboring value the wide-window search had locked onto recovered, in
every case, a value within $0.5\%$ of the seed (cells marked $^\dagger$
below use the recovered value). Two further cells, IP-11 at $n{=}20$
and IP-26 at $n{=}28$, sat just below the diagnostic's original 1\%
flag threshold and were identified only by a later review pass; their
tight-bracket re-polish recovered values within $0.03\%$ of the seed.
Separately, a follow-up diagnostic campaign established that the
branch-selection step feeding this table frequently does not reach the
labeled basis size: the routine orders candidate branches by
$(|\mathrm{Im}\,\xi|,|\xi|)$ and fills greedily until $n$ generalized
coordinates are reached, which for this geometry's branch inventory can
leave an even-valued shortfall unfilled. Forcing the routine to reach the
labeled $n$ anyway was tested against all 34 modes at both $n{=}28$ and
$n{=}36$, evaluated at each mode's own production frequency: in every
one of 32 such evaluations where this produced a larger basis than the
routine's own choice (29 at the modes' initial seed frequencies, and a
further 3 re-confirmed at their exact converged frequencies),
$\sigma_{\min}$ degraded from a genuine, deep minimum to a shallow
wide-window lock-on, confirming the shortfall is the routine correctly
declining to admit a matrix-corrupting branch, not a defect. At $n{=}28$, all eight
out-of-plane labels reach the full labeled basis; the in-plane labels
mostly reach 26 of 28, with IP-05/IP-06 reaching only 22. At $n{=}36$, no
label reaches the full labeled basis (out-of-plane: 28 or 32; in-plane:
30 or 34); the ``used'' columns below give the actually achieved size
per row. IP-06's own $n{=}36$ entry was recorded as 30 in one job and 34
in another, which an initial reading took for run-to-run irreproducibility.
It is not: the achieved size is a deterministic but rapidly varying
function of $\Omega$ in this neighborhood. Five repeat evaluations at the
tabulated frequency return 30 every time, and a scan across
$\Omega\in[1.196,1.205]$ returns 30 except inside a band roughly
$0.001$ wide near $\Omega\approx1.2032$--$1.2040$, where a tenth
transverse-wavenumber branch enters the search window and, being a
complex-quadrant pair worth four generalized coordinates, lifts the fill
from 30 to 34.\footnote{Jobs 2411851 and 2414031; the second covers
$[1.199,1.205]$, which the first did not reach. Transitions
$30\!\to\!34$ and back occur between consecutive $2\times10^{-4}$ scan
points, so the two earlier jobs' slightly different converged $\Omega$
placed them on opposite sides of one.} The two values are therefore both
correct, for marginally different frequencies; the range is retained in
the table rather than collapsed. Consequently the $\Delta_{28\to36}$
columns compare two sub-maximal bases in every row, and the
$\Delta_{20\to28}$ column for IP-05/IP-06 is a $20$-to-$22$, not
$20$-to-$28$, comparison (flagged individually below). The basis size
actually achieved by basis fill at each mode's own converged
$\Omega^{*}$, tabulated as the ``used'' columns below, is confirmed by
jobs 2324817/2324822.

\begin{table}[!htbp]\centering\small
\caption{Basis-size convergence, out-of-plane modes (raw $\Omega$;
labels OOP-01--OOP-08 correspond by row order to the main text's
Table 6, which does not print these labels). $^\dagger$: recovered via a
tight-bracket re-polish after an initial wide-window lock-on (see text
above). ``Used'' columns: the basis size actually achieved by
basis fill at the labeled target, confirmed at each mode's own
converged $\Omega^{*}$.}
\label{tab:conv-oop}
\begin{tabular}{cccccccc}
\toprule Label & $n{=}20$ & $n{=}28$ & $n{=}36$ & used@28 & used@36 & $\Delta_{20\to28}$(\%) & $\Delta_{28\to36}$(\%)\\ \midrule
OOP-01 & 0.383601 & 0.383598 & 0.383646 & 28 & 28 & $-0.0007$ & $+0.0125$\\
OOP-02 & 0.616726$^\dagger$ & 0.616724 & 0.616751 & 28 & 28 & $-0.0003$ & $+0.0044$\\
OOP-03 & 0.961242$^\dagger$ & 0.961210 & 0.961210 & 28 & 32 & $-0.0033$ & $+0.0000$\\
OOP-04 & 1.369616 & 1.369586 & 1.370489 & 28 & 28 & $-0.0022$ & $+0.0659$\\
OOP-05 & 1.752160 & 1.752037 & 1.755174 & 28 & 28 & $-0.0070$ & $+0.1791$\\
OOP-06 & 2.320817 & 2.320723 & 2.320714 & 28 & 32 & $-0.0041$ & $-0.0004$\\
OOP-07 & 2.379839 & 2.379688 & 2.379688 & 28 & 28 & $-0.0063$ & $+0.0000$\\
OOP-08 & 2.602593 & 2.602553 & 2.602553 & 28 & 28 & $-0.0015$ & $+0.0000$\\ \bottomrule
\end{tabular}
\end{table}

\begin{table}[!t]\centering\scriptsize
\setlength{\tabcolsep}{4pt}
\renewcommand{\arraystretch}{0.95}
\caption{Basis-size convergence, in-plane modes (raw $\Omega$; labels
IP-01--IP-26 correspond by row order to the main text's Table 7, which
does not print these labels). $^\dagger$: recovered via a
tight-bracket re-polish after an initial wide-window lock-on (see text
above). ``Used'' columns: as Table~\ref{tab:conv-oop}. $^\ddagger$: IP-06's
used@36 is 30 at the tabulated frequency and 34 inside a narrow adjacent
$\Omega$ band; both jobs are correct for their own converged $\Omega$
(see text). $^{\S}$:
$\Delta_{20\to28}$ computed on a 22-, not 28-coordinate basis: a
$20\to22$ comparison, not $20\to28$.}
\label{tab:conv-ip}
\begin{tabular}{cccccccc}
\toprule Label & $n{=}20$ & $n{=}28$ & $n{=}36$ & used@28 & used@36 & $\Delta_{20\to28}$(\%) & $\Delta_{28\to36}$(\%)\\ \midrule
IP-01 & 0.401200$^\dagger$ & 0.401185 & 0.401185 & 26 & 30 & $-0.0037$ & $+0.0000$\\
IP-02 & 0.687385 & 0.687377 & 0.687377 & 26 & 30 & $-0.0011$ & $+0.0000$\\
IP-03 & 0.764608 & 0.764588 & 0.764580 & 26 & 30 & $-0.0026$ & $-0.0010$\\
IP-04 & 1.142356$^\dagger$ & 1.143136$^\dagger$ & 1.142313 & 26 & 30 & $+0.0682$ & $-0.0720$\\
IP-05 & 1.145249 & 1.144393 & 1.145192 & 22 & 30 & $-0.0747^{\S}$ & $+0.0698$\\
IP-06 & 1.197860 & 1.197592$^\dagger$ & 1.203379 & 22 & 30--34$^\ddagger$ & $-0.0224^{\S}$ & $+0.4833$\\
IP-07 & 1.416707 & 1.416366 & 1.416355 & 26 & 34 & $-0.0241$ & $-0.0008$\\
IP-08 & 1.440521 & 1.440118 & 1.440110 & 26 & 34 & $-0.0280$ & $-0.0006$\\
IP-09 & 1.523138 & 1.522316 & 1.522258 & 26 & 34 & $-0.0540$ & $-0.0038$\\
IP-10 & 1.598919 & 1.598834 & 1.598826 & 26 & 34 & $-0.0053$ & $-0.0005$\\
IP-11 & 1.613504$^\dagger$ & 1.613132 & 1.613057 & 26 & 34 & $-0.0230$ & $-0.0046$\\
IP-12 & 1.690176 & 1.690020 & 1.689972 & 26 & 34 & $-0.0092$ & $-0.0029$\\
IP-13 & 1.779739$^\dagger$ & 1.779680 & 1.779676 & 26 & 34 & $-0.0033$ & $-0.0002$\\
IP-14 & 1.832690 & 1.831850 & 1.831788 & 26 & 34 & $-0.0458$ & $-0.0034$\\
IP-15 & 1.974997 & 1.972955 & 1.972920 & 26 & 34 & $-0.1034$ & $-0.0018$\\
IP-16 & 1.987420 & 1.987180 & 1.987158 & 26 & 34 & $-0.0121$ & $-0.0011$\\
IP-17 & 1.997673 & 1.996470 & 1.996243 & 26 & 34 & $-0.0602$ & $-0.0114$\\
IP-18 & 2.169765 & 2.166997 & 2.166775 & 26 & 34 & $-0.1276$ & $-0.0102$\\
IP-19 & 2.294238$^\dagger$ & 2.293218$^\dagger$ & 2.293042 & 26 & 34 & $-0.0445$ & $-0.0077$\\
IP-20 & 2.302741 & 2.301518 & 2.304771 & 26 & 34 & $-0.0531$ & $+0.1413$\\
IP-21 & 2.389179 & 2.388598 & 2.388507 & 26 & 34 & $-0.0243$ & $-0.0038$\\
IP-22 & 2.456152 & 2.450795 & 2.450496 & 26 & 34 & $-0.2181$ & $-0.0122$\\
IP-23 & 2.524864 & 2.521487 & 2.521288 & 26 & 34 & $-0.1337$ & $-0.0079$\\
IP-24 & 2.570691 & 2.566361 & 2.565772 & 26 & 34 & $-0.1684$ & $-0.0230$\\
IP-25 & 2.679963 & 2.676029 & 2.674815 & 26 & 34 & $-0.1468$ & $-0.0454$\\
IP-26 & 2.777863 & 2.772007$^\dagger$ & 2.771646$^\dagger$ & 26 & 34 & $-0.2108$ & $-0.0130$\\ \bottomrule
\end{tabular}
\renewcommand{\arraystretch}{1}
\end{table}

\paragraph{A non-default split/anchor enhancement.}
Main-text \S4 notes that a fully-default detector pass missed 4 of 26
tabulated extensional modes, later recovered as genuine roots of the
same $n=28$ determinant. A candidate enhancement to the blind pipeline
itself, an opt-in local split/anchor/adaptive-bisection strategy that
re-examines a coarse candidate's own local zoom for more than one
distinct minimum before committing to a single bracket, was tested in
the same whole-range, no-seeding configuration and recovered all four of
those modes unaided, including the closest pair (${\sim}0.0043$ apart in
$\Omega$). The same run left a different two-mode pair unrecovered under
the 0.1\% match tolerance (${\sim}0.0011$ apart, tighter than the
fine-zoom step itself); a targeted follow-up dense scan subsequently
confirmed both are genuine, distinct roots. The enhancement therefore
recovered 24 of 26 tabulated modes. It is not the package default: a
required second-geometry blind pass (in-plane orthotropic) recovered
only 1 of 8 known targets, and escalating the search basis left the
same mixed pattern, so a basis-size difference is not the cause.
Job-level detail: \S\ref{sec:provenance}.

\subsection{Instrument-by-instrument adjudication of the near-degenerate
ODD pair}\label{sec:instruments}
Main-text \S6.3 states the outcome of the eight-instrument suite applied
to the two candidates (raw $\Omega=2.828013$ and $2.838515$) against the
single FE target at $\Omega_{\mathrm{lit}}=111.46$. The individual
instruments and their outcomes are:

\begin{enumerate}\itemsep2pt
\item \emph{Basis-size bump} ($n_{\mathrm{dofs}}$ $20\to28$): both
candidates persist, essentially unchanged.
\item \emph{Sensitivity map}: a fixed-offset scan around each candidate's
golden-section result found no evidence that either was locking onto a
nearby artifact rather than a true root.
\item \emph{Null-vector gap analysis}: no informative separation.
\item \emph{Parity readout}: at the calibrated basis sizes both read cleanly
ODD, self-parity correlation $1.0000$, matching the FE target.
\item \emph{ODD-block topology scan}: no structural difference between the
two.
\item \emph{Modal assurance criterion}: the decisive instrument, and the
only one that compares mode \emph{shape} rather than frequency or matrix
structure. MAC against the FE eigenvector is $0.9913$ for $2.828013$ and
$0.9999$ for $2.838515$, a gap of $0.0086$ against the $0.2$ forced-call
threshold. Both candidates sit in the ``good correlation'' range of
experimental modal practice \cite{allemang2003} ($\mathrm{MAC}>0.9$; its
uncorrelated band is $<0.05$), so the absolute-level convention does not
separate them either, and the observed gap is roughly $20\times$ too small
to force a call at the gap threshold used here.
\item \emph{Pure-slope (moment-slope) residual}, applied independently after
the first six were exhausted: the two candidates' residual ratios differ by
$6.6\times$ ($2.549\times10^{-4}$ against $1.682\times10^{-3}$) and their
phase-aligned null-vector cosine is only $0.228$, indicating two radially
distinct modes rather than one being a numerical echo of the other.
\begin{sloppypar}
\item \emph{Enlarged-basis check} ($n_{\mathrm{ret}}=34$/$38$, beyond the
calibrated $n_{\mathrm{dofs}}=20$/$28$): candidate $2.828013$'s self-parity
reads ODD at $n_{\mathrm{ret}}=34$ and EVEN at $38$. Follow-up probes traced
this to a redistribution of null-vector energy onto a different dominant
root (from a pure-imaginary root carrying 31\% of the energy to a complex
root carrying 76\%) and ruled out a basis-selection artifact directly:
removing the single marginal root that crosses the basis-fill
boundary changes retained energy by under $0.1\%$ and leaves the parity
unchanged.
\end{sloppypar}
\end{enumerate}

\paragraph{MAC/FE nodal correspondence.} Instrument 6 above and the
two-sided extensional bar of main-text \S6.6 both compare
the exact solution's reconstructed mode shape against an FE eigenvector
by MAC, and both build the comparison the same way. The exact
reconstruction is not evaluated directly at the FE node coordinates; it
is first evaluated on a dense regular grid in the solver's own natural
coordinates ($r$ mapped affinely onto $[-\pi/2,\pi/2]$, $\theta$ over
$[-\Theta,\Theta]$, bisector at $\theta=0$). Each FE node's angular
coordinate (measured from the FE model's own origin at one
circumferential edge, not from the bisector) is shifted onto the
solver's bisector-centered convention before use. The grid's radial and
transverse displacements are then bilinearly interpolated onto every FE
node's coordinates; a node whose coordinates fall outside the grid's
covered domain returns a non-finite value and is excluded from that
particular MAC evaluation rather than extrapolated. The interpolated
radial/transverse pair is rotated into the FE model's own Cartesian
$(U_X,U_Y)$ using each node's own global angle, so the MAC correlation
of main-text \S6.6 is taken entirely in the FE model's native
frame.

A separate correspondence step is needed only to classify an FE
eigenvector's own symmetry class (SYMM/ANTI about the bisector), which
the two-sided bar's class-safe matching requires so a candidate is never
compared against an opposite-parity FE target. This pairs each FE node
with its own mirror-image node across the bisector, within the same
mesh, and is independent of the exact-solution correspondence above. At
the half-annulus ($2\Theta=\pi$, the \S6.5 geometry) this pairing is exact:
a node's reflected polar coordinates coincide with another printed
node's coordinates to within rounding. At every other sector angle,
including the primary quarter-annulus $2\Theta=\pi/2$ and the rest of
the main-text \S6.4 sweep, the FE decks are meshed in Cartesian
coordinates, so a node's exact polar mirror image generally does not
land on another node's printed coordinates. The production fallback,
used at those angles, is a nearest-neighbor spatial search in Cartesian
coordinates for the closest node to the exact mirror image, accepted
only within half a typical element edge length; a node with no partner
inside that radius is dropped from the symmetry calculation rather than
mis-paired. The primary-geometry MAC calibration of main-text \S6.6
therefore used this fallback, not exact pairing. At the half-annulus,
where exact pairing succeeds, forcing the Cartesian fallback instead
left all 57 FE symmetry labels and all eight calibration MAC calls
unchanged at threshold $0.744$.

\paragraph{Owning-block scoring and held-out candidates.}
Two properties of the extensional reconstruction have to be respected
before MAC means anything. The parity blocks $q{=}0$ and $q{=}1$ of the
assembled $K$ are the sector's two exact symmetry classes, so a MAC
taken across them is analytically zero: a reading of $0.0000$ records a
block-selection outcome, not an absence of correlation. And the smallest
singular vector of a \emph{non}-singular block varies smoothly with
$\Omega$, so a block evaluated near a neighboring candidate's root
reproduces that neighbor's mode shape: one candidate at the primary
geometry reads $0.9658$ this way without owning a root in that block at
all. Each candidate must therefore be scored on the parity block that
owns a root at its own $\Omega$.

Selecting that block requires a rule, and no single criterion suffices.
Choosing the more singular block misclassifies one candidate at the
second geometry; choosing the smaller edge residual misclassifies a
different one at the first. Requiring that a block's
$\log_{10}\sigma_{\min}$ sit at a local minimum in $\Omega$ before it is
eligible, and breaking ties on the residual, classifies 24 of 24 against
23 of 24 for either criterion alone. At $(2.0,\,0.5)$ all three
selection criteria score 8 of 8, so that geometry confirms the bar
transfers but supplies no evidence that the rule is better than the
simpler alternatives; the rule's margin rests on two candidates.

Two candidates were held out of the calibration set entirely, because
the pointwise residual and the FE evidence disagree about them. At
$\Omega=1.667187$ the production residual verdict is ARTIFACT-like, the
independent FE frequency match is $0.011\%$, and MAC on the root-owning
block reads $0.9981$; at $\Omega=2.117877$ the corresponding figures are
$0.006\%$ and $0.9980$. Both sit far above any threshold in the window,
so MAC agrees with the FE evidence where the residual does not.
This hold-out pair is not the eight residual-unresolved candidates of
main-text \S6.5: those eight include $\Omega=1.667187$ as one of two
MAC-REAL calls (the other is $\Omega=1.894176$, MAC $0.9953$);
$\Omega=2.117877$ is a calibration hold-out only. Main-text \S6.6
records the production call (midpoint $0.744$ of the $\nu=0.30$
window); \S6.5 records the outcome on the eight.

The in-plane analogue of Table~\ref{tab:oop-spurious-sm}, listing
owning-block residuals for these eight calibration candidates, is
Table~\ref{tab:ip-spurious}.\footnote{Job 2434754. Owning-block residuals
match the calibration residuals to $10^{-6}$ relative.}

\paragraph{Calibration-window robustness.} Main-text \S6.6 quotes the
seed-based REAL floor and ARTIFACT ceiling ($0.9079$, $0.5801$) because they
require no extra refinement step. Evaluating the one calibration
candidate whose catalogued frequency lies furthest from its own
converged root at that root instead raises the combined REAL floor from
$0.9079$ to $0.9869$. The combined window (ARTIFACT ceiling $0.5801$ to
REAL floor) therefore widens from $0.3278$ to $0.4068$; the
per-geometry gap at $(2.0,1.0)$ on the seed-based table is already that
same $0.4068$ and is unchanged to the printed digits. The reported window is
otherwise insensitive to how the FE comparison target is chosen: MAC
taken against each candidate's own nearest-frequency FE target gives the
same window as MAC taken as the best correlation over every sampled
target, so the calibration does not depend on the nearest-frequency
pairing being correct. Unrestricted nearest-frequency matching is
nevertheless not class-safe on its own: at
$(r_0/2b,2\Theta/\pi)=(2.0,1.0)$ one calibration artifact's
nearest-frequency FE target belongs to the opposite symmetry class, and
MAC against that target is identically zero, the wrong-class
signature described above, not an absence of correlation. This is why
main-text \S6.6 scores every candidate as the maximum over FE modes of
the same symmetry class as the root-owning block, not simply the
nearest frequency.

\subsection{Cut-off-adjacent screening and the extensional basis-packing
shortfall}\label{sec:cutoff-appendix}
Main-text \S6.4 sets aside a cut-off-adjacent subset of the 24-geometry
(48 geometry/motion) sweep and reports only its final tallies. That subset was screened with
the same \S5 pointwise discriminator used for the primary tables, its
first application outside the single geometry at which it was calibrated.
All 119 flexural candidates
were resolved to their target basis and adjudicated: 57 REAL-like, 62
ARTIFACT-like, 0 AMBIGUOUS (a 47.9\% REAL-like fraction; these 119 are a
subset of the 695 flexural candidates screened in main-text \S6.4, whose
complement reads 272/535, or 50.8\%).\footnote{Direct
pass: jobs 2416866--2416869. Escalation ladder for the 12 that did not
reconstruct on a cold search: job 2417234.} The REAL-like and
ARTIFACT-like populations do not overlap ($r_M\le0.041$ against
$r_M\ge2.05$ for the directly reconstructed majority, a $\ge50$-fold
separation, narrower than, but qualitatively identical to, the
primary geometry's own ${\sim}2000$-fold gap).

The extensional side has no calibrated \emph{residual} verdict bar at any
geometry, so its figures are raw residual ratios. All 96 extensional
candidates were resolved: 88 bimodal, 8 middle-band by the same
$r_{Tyy}<0.3$ / $>4$ split used at the primary geometry ($r_{Tyy}$ and
$r_{Tyr}$ here are the main text's extensional residual ratios
$r_{\theta\theta}$ and $r_{\theta r}$, in the solver's own edge-frame
naming).
Seventy-one reached their target basis (direct search plus an escalation
ladder); the remaining 25 were stuck at exactly $n_{\mathrm{dofs}}-2$
even at a guarded $x_{\max}$ cap of 90 and were evaluated at that
achieved, smaller basis, following the same used@28/used@36 relabeling
precedent as \S\ref{sec:conv-appendix} and the IP-06 footnote, with one
exception: a candidate that reached the labeled fill with a shallow
$\sigma_{\min}$ ($-2.681$) and is reported at that fill.\footnote{Job
2419006.}

The 25-candidate shortfall comes from a packing artifact of greedy branch selection; the underlying inventory is not shrinking. Each admitted root contributes
either 2 degrees of freedom (real or pure-imaginary) or 4 (a
complex-quadrant pair); the routine fills by ascending
$(|\mathrm{Im}\,\zeta|,|\zeta|)$ and stops when the next candidate would
overshoot the target. At each shortfall point the running fill reaches
exactly $n_{\mathrm{dofs}}-2$ with only 4-dof pairs remaining. The same
cap arises at a control geometry where 30--34 degrees of freedom are
available, so a genuine geometry-dependent drop in complex-branch
inventory with increasing $r_0/2b$ coexists but is not the cause. No
adjudicated (flexural) conclusion depends on this. The full-sweep match
recapture behind Table~\ref{tab:geomsweep} is job 2416746 (superseding an
intermediate figure from an earlier recapture of the same sweep, 1231 of
1474, 83.5\%, mean 0.578\%, worst 3.000\%, which had itself already
superseded the original interactively-run, never-committed statistic,
1232 of 1466, 84.0\%; the per-candidate set is not identical across any
of these checkpoints, and the aggregate figure is what each recapture
re-derives). That job also settles how much the one unrecoverable
methodological choice matters: rerun at extensional-contamination
tolerances of $0.5\%$, $1.0\%$ and $2.0\%$, it returns $1261$, $1256$ and
$1253$ matches ($85.5\%$, $85.2\%$, $84.9\%$), with mean-of-per-geometry-means
$0.573\%$/$0.574\%$/$0.574\%$ and worst case $2.990\%$ in all three --- a
fourfold change in the unrecorded parameter moves the headline by $0.6$
percentage points. The reported figures are the $1.0\%$ column. The flexural
residual screen behind its reconstruction/verdict tallies is jobs
2428355--2428358.

\begin{table}[!htbp]\centering\small
\caption{Free-free detector accuracy vs.\ an independent FE benchmark,
4 radius ratios $\times$ 6 sector angles (both motion types), excluding
each geometry's cut-off-adjacent candidates; main-text \S6.4.
$n$: matches within 3\% of 1475 attempted. The sweep produces 2291
candidates in all (695 flexural, 1596 extensional); 1475 remain after
removing each geometry's cut-off-adjacent candidates (119 flexural, 96
extensional) and those lying beyond the FE extraction range (12 flexural,
589 extensional). Every accepted flexural
candidate was residual-screened: 654 of 695 reconstructed, 329 REAL-like /
321 ARTIFACT-like / 4 AMBIGUOUS. Cut-off-adjacent subset and
extensional packing shortfall: \S\ref{sec:cutoff-appendix}.}
\label{tab:geomsweep}
\begin{tabular}{ccccc}
\toprule $r_0/2b$ & $R_i/R_o$ & $n$ & mean err (\%) & max err (\%)\\ \midrule
1.5 & 0.500 & 274 & 0.572 & 2.976\\
1.66667 & 0.538 & 294 & 0.559 & 2.944\\
2.0 & 0.600 & 336 & 0.576 & 2.965\\
2.5 & 0.667 & 352 & 0.590 & 2.990\\ \midrule
all & --- & 1256/1475 & 0.574 & 2.990\\
\bottomrule
\end{tabular}
\end{table}

The $\nu=0.35$ flexural candidate $\Omega=1.432743$ in
Table~\ref{tab:r150a100} below, which job 2347329 could not reconstruct
at $n{=}16$, was recovered by the cut-off-adjacent escalation ladder
(job 2421415) as ARTIFACT-like and FE-matched within $3\%$; the table's
IP FE-matched rates use the 80-mode extract (jobs 2421055/2421056).

\begin{table}[!htbp]\centering\small
\caption{Second-geometry adjudication summary ($r_0/2b=1.5$,
$2\Theta/\pi=1.0$), both motion types, both materials; main-text \S6.5.
FE-matched: within 3\% of an independent FE frequency. The $\nu=0.35$
flexural candidate $\Omega=1.432743$ was recovered by the
cut-off-adjacent escalation ladder as ARTIFACT-like and FE-matched
within $3\%$; it is included. IP FE-matched rates use the 80-mode
extract.}
\label{tab:r150a100}
\begin{tabular}{clccccc}
\toprule $\nu$ & motion & $N$ & REAL-like & ARTIFACT-like & AMBIG. & FE-matched\\ \midrule
0.35 & flexural (OOP)    & 28 & 14 & 14 & 0  & 21/28 (75.0\%)\\
0.35 & extensional (IP)$^\ast$  & 61 & 31 & 19 & 11 & 57/61 (93.4\%)\\
0.30 & flexural (OOP)    & 30 & 15 & 15 & 0  & 20/30 (66.7\%)\\
0.30 & extensional (IP)$^\ast$  & 64 & 28 & 18 & 18 & 61/64 (95.3\%)\\
\bottomrule
\multicolumn{7}{l}{\footnotesize $^\ast$raw residual-ratio classification, not a calibrated residual verdict (main-text \S6.3; mode-shape bar: \S6.6).}\\
\end{tabular}
\end{table}

\begin{table}[!htbp]\centering\small
\caption{Second radius-ratio adjudication summary ($r_0/2b=2.0$,
$R_i/R_o=0.600$, $2\Theta/\pi=1.0$, $\nu=0.30$); main-text \S6.5.
FE-matched: within 3\% of an independent FE frequency.}
\label{tab:r200a100}
\begin{tabular}{clccccc}
\toprule $\nu$ & motion & $N$ & REAL-like & ARTIFACT-like & AMBIG. & FE-matched\\ \midrule
0.30 & flexural (OOP)    & 33 & 18 & 15 & 0  & 25/33 (75.8\%)\\
0.30 & extensional (IP)$^\ast$  & 75 & 35 & 18 & 22 & 68/75 (90.7\%)\\
\bottomrule
\multicolumn{7}{l}{\footnotesize $^\ast$raw residual-ratio classification, not a calibrated residual verdict (main-text \S6.3; mode-shape bar: \S6.6).}\\
\end{tabular}
\end{table}

\begin{table}[!htbp]\centering\small
\caption{Crossed-geometry adjudication summary ($r_0/2b=2.0$,
$R_i/R_o=0.600$, $2\Theta/\pi=0.5$, $\nu=0.30$); main-text \S6.5.
FE-matched: within 3\% of an independent FE frequency. All 10 flexural
REAL-like candidates are among the 11 FE-matched;
the eleventh is the ARTIFACT-like FE-mode-16 companion discussed in
main-text \S6.5.}
\label{tab:r200a50}
\begin{tabular}{clccccc}
\toprule $\nu$ & motion & $N$ & REAL-like & ARTIFACT-like & AMBIG. & FE-matched\\ \midrule
0.30 & flexural (OOP)    & 18 & 10 & 8 & 0  & 11/18 (61.1\%)\\
0.30 & extensional (IP)$^\ast$  & 47 & 16 & 16 & 15 & 43/47 (91.5\%)\\
\bottomrule
\multicolumn{7}{l}{\footnotesize $^\ast$raw residual-ratio classification, not a calibrated residual verdict (main-text \S6.3; mode-shape bar: \S6.6).}\\
\end{tabular}
\end{table}

\begin{table}[!htbp]\centering\small
\caption{Owning-block traction residuals at $n{=}20$ for the
primary-geometry MAC calibration set of main-text \S6.6 (the listing
analogue of Table~\ref{tab:oop-spurious-sm}). Labels are that calibration assignment,
not a row-by-row application of the \S5 ${\sim}0.3$/${\sim}4$ split.
The REAL-like half are main-text Table 7 rows 1, 4, 10 and 23; the ARTIFACT-like
half are catalogued spurious zeros of the same determinant. Residuals
are on the root-owning parity block (local-minimum residual rule);
the last numeric column is $\max(r_{Tyy},r_{Tyr})$.}
\label{tab:ip-spurious}
\begin{tabular}{lcccccc}
\toprule Label & $\Omega$ & block & $r_{Tyy}$ & $r_{Tyr}$ & $\max(r_{Tyy},r_{Tyr})$ & Verdict\\ \midrule
real\_a & 0.401154 & $q{=}1$ & 0.03052 & 0.01808 & 0.03052 & REAL-like\\
real\_b & 1.142287 & $q{=}0$ & 0.04982 & 0.03534 & 0.04982 & REAL-like\\
real\_c & 1.598985 & $q{=}1$ & 0.04855 & 0.02166 & 0.04855 & REAL-like\\
real\_d & 2.521489 & $q{=}1$ & 0.6858 & 0.3205 & 0.6858 & REAL-like\\
art\_a & 0.542105 & $q{=}0$ & 6.907 & 4.801 & 6.907 & ARTIFACT-like\\
art\_b & 1.021144 & $q{=}1$ & 6.058 & 8.287 & 8.287 & ARTIFACT-like\\
art\_c & 1.454616 & $q{=}0$ & 7.035 & 3.288 & 7.035 & ARTIFACT-like\\
art\_d & 1.603259 & $q{=}0$ & 5.199 & 0.917 & 5.199 & ARTIFACT-like\\
\bottomrule
\end{tabular}
\end{table}

\begin{table}[!htbp]\centering\small
\caption{FF-P1 flexural spurious zeros ($n_{\mathrm{dofs}}=20$),
the listing referred to from main-text \S6.3. Depth is
$\log_{10}\sigma_{\min}$ at the same basis. Physical modes at this
basis have $r_M$ spanning $0.0025$--$0.037$ and depths $-2.81$ to
$-4.78$ (main-text Table 6). The last row, raw $\Omega=2.817613$, is
main-text Fig.~2(e).}
\label{tab:oop-spurious-sm}
\begin{tabular}{cccccc}
\toprule Raw $\Omega$ & Parity & $r_V$ & $r_M$ & depth & Verdict\\ \midrule
0.320670 & EVEN & 53.4 & 74.9 & $-3.69$ & ARTIFACT-like\\
0.844488 & ODD & 75.2 & 76.0 & $-4.25$ & ARTIFACT-like\\
1.001419 & EVEN & 454 & 290 & $-5.29$ & ARTIFACT-like\\
1.602090 & EVEN & 95.6 & 75.5 & $-4.71$ & ARTIFACT-like\\
2.121685 & ODD & 620 & 348 & $-4.13$ & ARTIFACT-like\\
2.817613 & EVEN & 180.6 & 75.52 & $-4.82$ & ARTIFACT-like\\
\bottomrule
\end{tabular}
\end{table}
\FloatBarrier

\subsection{Extensional close-pair resolutions}\label{sec:closepairs}
Two extensional benchmark pairs are close enough in frequency that a
single root is asked to cover both FE or literature targets, for
opposite reasons. Main-text \S6.3 and \S6.9 state only the outcome; both
are resolved here in full.

\paragraph{The 1114.9/1115.3~Hz pair (main-text \S6.3, Table 7).}
Isolated golden-section refinement at the two predicted pair members, in
brackets that cannot reach each other, finds two shift-stable deep
minima: $\Omega^{*}=2.772422$ ($1115.18$~Hz, depth $-4.43$) and
$\Omega^{*}=2.773399$ ($1115.58$~Hz, depth $-4.82$), each $0.015\%$ from
one FE member. The production detector still returns a single
singularity ($2.772007$ at $n{=}28$; $2.772924$ at $n{=}36$, between the
two isolates). MAC against the two isolates assigns both REAL-like to
the same FE mode (FE\#29, MAC $0.9898$ and $0.9901$); the production
root $\Omega=2.771955$ is that same shape (MAC $0.9896$). One shape, two
frequencies; the main-text table is unchanged.

\paragraph{Qin's 4.4293/4.4754 pair (main-text \S6.9, Table 10).}
This is not Table 7's FE-matched pair (rows 4--5, $\Omega=1.142287$ and
$1.145158$, $0.25\%$ apart at the nominal wavenumber window). Those two
roots are the comparison against FE in main-text \S6.3; they are not an
assignment of Qin's published values, and no claim is made about which
spacing is correct. The $n_{\mathrm{dofs}}=12$ search lists a single
nearby root (1.147615) against both published Qin targets. Repeating
that targeted search at requested $n{=}20$ and $n{=}28$ --- an enlarged
$x_{\max}$ search that realizes 18 degrees of freedom whichever of the
two is requested --- still returns one singularity
($\Omega=1.14274$, $\Omega_{\mathrm{Qin}}=4.4529$), $0.53\%$ and
$0.50\%$ from the two Qin targets. A second isolated dip sits $0.42\%$ away
at $\Omega=1.14758$ ($\Omega_{\mathrm{Qin}}=4.4717$, depth $-5.18$ at
$n{=}20$ and $-6.07$ at $n{=}28$, shift $0.0001\%$); the Qin-targeted
search still returns only the $1.14274$ root, so Qin's pair remains
unseparated. Isolation brackets of $\pm0.15\%$ cannot reach the
neighboring root; this is not a table split, the published row stays
the $n{=}12$ shared listing. MAC on the two isolates assigns both
REAL-like to distinct FE modes of opposite parity: $\Omega=1.14274$
ANTI, MAC $0.9996$ at FE\#7 (Table 7's $459.5$~Hz partner; the isolate
itself converts to $459.7$~Hz); $\Omega=1.14758$ SYMM, MAC
$0.9871$ at FE\#8 (Table 7's $460.6$~Hz partner; the isolate itself
converts to $461.6$~Hz). Here the outcome runs the other way:
MAC finds two distinct shapes, but the Qin-targeted search still
returns only the one root.

The two pairs illustrate the same limit from opposite sides: one
production root can legitimately stand for two nearby targets when
their mode shapes coincide (the 1114.9/1115.3~Hz pair), or, when they do
not (Qin's pair), a targeted search still cannot assign a distinct
production root to each published Qin value even though MAC shows the
two isolates are physically distinct shapes.

\subsection{Shi et al.\ out-of-plane refinement detail}\label{sec:shi2014-detail}
Main-text \S6.10 states only the final refined roots
(Table 11); two of Shi et al.'s six targets needed
more than the basic window-scan-plus-golden-section pipeline, detailed
here. The acceptance criteria applied throughout both cases are
$\log_{10}\sigma_{\min}\le-3.5$, error below 2\%, and both bracket-shift
replicas within $0.1\%$ of the converged root.

\paragraph{53.649: a genuine near-degenerate pair.} The initial
$\pm2.5\%$-bracket golden-section refinement failed its bracket-shift
acceptance test (the refined position moved $0.59\%$ when the bracket
was nudged), a symptom main-text \S5's instruments treat as lock-on instead of a located root. An unconstrained dense scan of the neighborhood
isolated two genuine, stable roots $0.59\%$ apart: $\Omega^{*}=1.957111$
($\Omega_{\mathrm{lit}}=53.6553$, err $0.012\%$, depth $-5.848$) at the
published frequency, and $\Omega^{*}=1.945645$ ($\Omega_{\mathrm{lit}}=53.341$,
err $0.574\%$, depth $-3.861$), the wide bracket's lock-on target. Each
root is stable to $0.0004\%$ or better inside an isolated $\pm0.4\%$
bracket that cannot reach the other, confirming the original $0.59\%$
walk was the separation of two roots, not an unlocated singularity. The
published value is reported in Table 11 against the
deeper, closer member; the second root stands on its own as a genuine singularity, left unforced to the published value.

\paragraph{38.428: a basis-depth near-miss, not a null result.} At the
production screening basis $n{=}20$, a dip sits exactly on the
predicted frequency, 1.9 decades below the local background (deepest
point $\log_{10}\sigma_{\min}=-3.297$ at raw $\Omega=1.40170$), but
falls short of the $-3.5$ acceptance depth. A 33-point scan of the
tight $\pm5\%$ window confirmed this is a genuine depth shortfall, not
a scan-resolution miss (the unrelated feature 12.5\% away that won the
original wide window's arg-max lies outside this tighter window).
Repeating at $n{=}28$/$36$ clears the acceptance bar ($-3.823$,
$-4.312$); golden-section refinement then gives a shift-stable root at
$\Omega_{\mathrm{lit}}=38.4295$ (err $0.004\%$, depth $-4.706$ at
$n{=}28$; $n{=}36$ agrees to $38.4289$, depth $-5.166$), reported in the
table as a refined root.

\subsection{Raw-to-physical conversion constants}\label{sec:conversion-sm}
The isotropic free-free benchmark underlying main-text Tables 6 and 7
(both the SHELL281 flexural and the PLANE183 extensional FE models) is a
single physical plate: inner radius $R_i=4$~m, outer radius $R_o=8$~m
(so $R_i/R_o=0.5$ and $r_0/2b=1.5$), thickness $H=0.08$~m, sector angle
$2\Theta=\pi/2$ (half-angle $\Theta=\pi/4$, the quarter-annulus of
main-text \S6.2--\S6.3), steel material $E=210$~GPa, $\nu=0.30$,
$\rho=7800$~kg/m$^3$. The governing slenderness is thickness over the
smallest in-plane dimension, the radial width $2b=4$~m, giving $H/2b=0.02$
($H/R_o=0.01$) --- well inside the range where the classical thin-plate
model applies, though see main-text \S6.3 on the residual theory difference
against a shear-deformable finite-element benchmark. With $b=(R_o-R_i)/2=2$~m, $c_{66}=G=E/2(1+\nu)$ and
$\bar c_{11}=E/(1-\nu^2)$, the two nondimensionalizing constants
entering main-text \S3.1's $\bar\Omega=\omega/\bar\omega$
(Ref.~\cite{seok1} Eqs.~(26)--(28); Ref.~\cite{seok2} Eqs.~(23)--(25))
are
\[
\bar\omega=\frac{\pi}{2b}\sqrt{\frac{c_{66}}{\rho}}=2527.351\ \mathrm{rad/s},
\qquad
\bar k=\frac{2b}{\pi}\sqrt{\frac{3c_{66}}{h^{2}\bar c_{11}}}=32.617,
\qquad h=\frac{H}{2}.
\]
The raw solver output converts to physical Hertz directly for the
in-plane (extensional) tables, $f=\bar\Omega\bar\omega/2\pi$, verified
against the FEM entries in main-text Table 7, confirming the in-plane
raw column already \emph{is} $\bar\Omega=\omega/\bar\omega$ with no
further conversion needed. The out-of-plane raw column is \emph{not}
$\bar\Omega$: $\bar\Omega$ enters the flexural dimensionless equation
only through the product $\bar k\bar\Omega$, which is what
Ref.~\cite{seok1}'s Table~2 and main-text Table 6 both report.
Recovering physical frequency from the flexural raw column therefore
carries one extra division by $\bar k$,
\[
f_{\mathrm{OOP}}=\frac{(\bar k\bar\Omega)\,\bar\omega}{2\pi\,\bar k},
\]
and this is what populates Table 6's $f$(Hz) column, giving both
flexural and extensional tables a shared physical-frequency metric even
though their raw columns are different quantities. Omitting the $\bar k$
here would inflate every flexural frequency by a factor of $32.617$ for
this plate and would place the first flexural mode above the first
extensional one, which is not physical for a plate of these proportions.
The isotropic out-of-plane table (main-text Table 6) is, in its
principal frequency column, reported in the classical literature
convention
$\Omega_{\mathrm{lit}}=\omega R_o^2\sqrt{\rho H/D}$
($D=EH^3/12(1-\nu^2)$, the convention Shi, Liang, Wang and Teng and the
original Seok--Tiersten tables both use), which is linear in the raw
column $\bar k\bar\Omega$ through the same factor $\bar k$:
\[
\Omega_{\mathrm{lit}} \;=\; (\bar k\bar\Omega)\,\frac{\bar\omega}{\bar k}\,R_o^2\sqrt{\frac{\rho H}{D}}
\;=\; \omega\,R_o^2\sqrt{\frac{\rho H}{D}}.
\]
Table 5 uses the same $\Omega_{\mathrm{lit}}$ symbol under the
orthotropic constants of main-text \S3.4, not this isotropic $\bar k$
conversion.
The conversion is cross-checked numerically against all 8 rows of
main-text Table 6 (e.g.\ raw $0.383695\to\Omega_{\mathrm{lit}}=15.148$,
raw $0.961312\to37.951$, raw $2.602500\to102.743$, each matching the
tabulated solver value to the last printed digit). The extra factor of
$\bar k$ for the out-of-plane case (absent from the in-plane conversion)
reflects that $\bar\omega$ alone nondimensionalizes the extensional
(second-order, wave-speed-scaled) problem, while the flexural
(fourth-order, thickness-scaled) problem's frequency parameter
additionally carries the thickness-dependent wavenumber scale $\bar k$;
both constants derive from the same base quantities $b,c_{66},\rho$
(and, for $\bar k$, $h,\bar c_{11}$), so no new physical input is
required beyond the geometry and material stated above.

\subsection{SUBDOM: numerical procedure and gate detail}\label{sec:subdom-detail}
This section gives the computational detail behind main-text \S6.7 at a
level a reimplementation could follow; the main text states only the
construction and the results.

\paragraph{Factorization and whitening.} For a fixed $\Omega$, the
in-plane free-free block assembly is $\mathbf{K}=\mathbf{D}^{\mathsf T}\mathbf{T}$
(bold denotes the matrices of this construction, distinguishing
$\mathbf{D},\mathbf{T},\mathbf{M},\mathbf{R}$ from the scalars $D$, $T$,
$M_{r\theta}$ and $R^{*}$ used elsewhere), verified
entrywise against the deployed \texttt{\_build\_K\_real} assembly on the
same basis (gate G1, bar $10^{-30}$, measured $1.21\times10^{-41}$ to
$8.80\times10^{-41}$ across all geometries reported). $\mathbf{D}$ is
QR-decomposed, $\mathbf{D}=\mathbf{Q}\mathbf{R}$; the pencil is whitened by
forming $\mathbf{M}=\mathbf{T}\mathbf{R}^{-1}$; the
singular value decomposition of $\mathbf{M}$ gives generalized singular values
$\sigma_1\le\sigma_2\le\cdots\le\sigma_n$, of which $\sigma_1=\eta(\Omega)$
and the associated right singular vector, mapped back through $\mathbf{R}^{-1}$,
is the traction-minimizing coefficient vector.

\paragraph{Degeneracy pooling and interior projection.} With $\gamma=4$
(adopted from the labelled-set window $\gamma=1.5$--$8$ measured before any
blind run; $\gamma=3.0$ was the original pre-registered value),
$k=\#\{i:\sigma_i\le G\sigma_1\}$ selects the
near-degenerate cluster of lowest-traction directions. Each of the $k$
associated coefficient vectors, together with the null vector
$c_{\mathrm{null}}$ of the accepted zero, is mapped to an interior
displacement field via \texttt{field\_of()}/
\texttt{adv.build\_grid\_from\_cfull} rather than the edge trace (an
edge-trace version is defeated by the mode-2 bleed-through artifact,
main-text \S6.7). The $k$ interior fields are Gram--Schmidt
orthonormalized under the same identity-weighted inner product used for
the main-text \S6.6 MAC, on the interior $r\times\theta$ grid those
helpers produce; SUBDOM is the cumulative squared projection of the
normalized $c_{\mathrm{null}}$ interior field onto that orthonormal span.

\paragraph{Gates, restated with bars.} G1: entrywise
$\mathbf{K}=\mathbf{D}^{\mathsf T}\mathbf{T}$, bar $10^{-30}$. G2: reproduces the published FF-P1 extensional residual
ratios, bar $10^{-4}$; applicable only at the primary geometry and
production basis ($n_{\mathrm{dofs}}=20$, $x_{\max}=20$), reported
not-applicable elsewhere by design. G3: the algebraic identity
$\eta(\Omega)\le\eta(c_{\mathrm{null}})$, which must hold exactly since
$c_{\mathrm{null}}$ is one feasible point of the minimization defining
$\eta$. G4: cross-checks the edge-trace decomposition SUB1 (not SUBDOM
itself) against an explicit MAC computation, bar $10^{-9}$; SUB1 and
SUBDOM are expected to disagree in general, since one is evaluated on the
edge trace and the other on the interior field. G5: confirms the
candidate's $\Omega$ is a genuine local minimum of $\sigma_{\min}$
under this basis, rejecting candidates that moved under the current
search (reported \texttt{C\_MOVED}, not a failure of the criterion). G6:
negative-control siting (known non-root points must fail G5 and must
also read SUBDOM below threshold); not applicable in blind-mode runs,
where no candidate is pre-labelled.

\paragraph{Calibration set.} The threshold, SUBDOM~$>0.70$, was set at
$n_{\mathrm{dofs}}=20$ as the midpoint of the window spanned by 52
labelled points across four
(geometry, $\nu$) combinations: the primary geometry, the two further
$\nu=0.30$ geometries of main-text \S6.5, and the $\nu=0.35$ geometry of
main-text \S6.4. The 52 points
include the two historical near-degenerate ``killer'' modes (one of
which, at an enlarged basis $n_{\mathrm{dofs}}=28$/$x_{\max}=60$, is the
sole exception recorded against this criterion to date, characterized and
not hidden) and 13 non-root negative controls (points where $K(\Omega)$
has no accepted zero at all). All 52 classify correctly at that
cut. It was then frozen: $0.70$ was fixed before any of the cluster runs
in main-text \S6.7 were submitted --- the four-geometry curated
confirmation and the 122-candidate blind set alike --- and was not
retuned afterward. ``Pre-registered'' below refers to that freeze, not to
a value chosen before the calibration itself. $\gamma=4$ was adopted after a second-geometry blind set showed
$\gamma=2$--$6$ inert and after the sole $\gamma=3.0$ knife-edge (FE mode 23,
$\gamma=3.01$) had been identified; it is a value inside the labelled-set
window. A re-score of the stored blind
checkpoints at $\gamma=4$ is $45/49$ and $25/28$ with zero duplicates
accepted (at $\gamma=3.0$, $44/49$ and $25/28$).

\paragraph{Block selection at near-degenerate pairs.} Where two solver
structures sit within $\lesssim10^{-3}$ in $\Omega$ (as at
$r_0/2b=2.0$, $2\Theta/\pi=1.0$'s FE mode 35, main-text \S6.7), the
production pipeline's \texttt{RULE\_LOCALMIN\_RESID} rule selects which
parity block to reconstruct by which is locally more singular ---
inherited from main-text \S6.6's own pipeline, not independently
re-derived for SUBDOM. This is the same style of automatic
sigma-based selection whose unmodified reuse, in the standalone
diagnostic probe built to cross-check the FE mode 35 pairing, silently
reconstructed the same block for both candidates on its first attempt (an
exact $\mathrm{MAC}=0.0000$/$0.0000$ result for both, later traced to the
bug and not reported as a finding). The fix there --- an explicit
\texttt{force\_block} parameter with an in-run fidelity gate proving
$\texttt{force\_block=None}$ reproduces the original automatic selection
bit-for-bit --- is specific to that standalone probe; the production
pipeline's own \texttt{RULE\_LOCALMIN\_RESID} rule is unmodified. Forcing
both blocks at the remaining disputed modes (main-text \S6.7) shows that
this inherited rule, not the interior projection, is what misses FE
modes 32 and 26 at $n_{\mathrm{dofs}}=20$, $x_{\max}=20$: the
auto-picked block reads $\mathrm{MAC}=0.0000$,
while the other block at the same $\Omega$ matches the frequency-nearest
FE mode ($\mathrm{MAC}=0.8528$ and $0.9302$) and, at FE mode 32,
SUBDOM~$=0.7136$ on that matching block, above the pre-registered cut
(FE mode 26's matching block reads SUBDOM~$=0.8728$). FE mode 26
recovers at an enlarged basis, below (main text \S6.7); FE mode 32
does not. FE mode 50 had no
owner above $0.744$ in FE modes $4$--$60$ under this first pass (best MAC
$0.4963$) and neither block certified (SUBDOM~$=0.0894$/$0.4813$); an
$\Omega$ scan and golden-section refinement subsequently located a
second, previously uncatalogued structure at $\Omega=2.144291$
(SUBDOM~$=0.9974$, $\mathrm{MAC}=0.9993$) that resolves this pairing
(main text \S6.7). FE mode 30's
auto-picked block is a pairing artifact; the other block reads
SUBDOM~$=0.7942$ but MAC against FE\#30 and against the eleven existing
r200/a50 eigenvectors is below $0.3$ --- max-over-blocks would accept
it, and is not adopted.

\paragraph{A fourth geometry's MAC cross-check.} The $\nu=0.35$
geometry ($r_0/2b=1.5$, $2\Theta/\pi=1.0$) in the main-text SUBDOM
calibration table had SUBDOM values only when that table was first
populated: no ANSYS eigenvector extract existed for it at the time.
An 80-mode PLANE183 extract (job 2433080, modes $4$--$60$ exported as
mesh96 eigenvectors) has since made the same SUBDOM-vs-MAC bake-off
used at the other three geometries (main-text §6.6's
calibration; jobs 2447412--2447414) available here too. Run in the
diagnostic container against all 12 of this geometry's labelled
points ($n_{\mathrm{dofs}}=20$, $x_{\max}=14$, matching the SUBDOM
table), with $\mathrm{MAC\_XMAX}=14$ set explicitly before loading the
reconstruction helper (the module-load-time default-binding pitfall
documented above for \texttt{XMAX0}): all 6 REAL-labelled points read
$\mathrm{MAC}>0.999$ on the SUBDOM-selected block; all 6
ARTIFACT-labelled points reject on that block, three below
$\mathrm{MAC}=0.3$ and three via $\mathrm{SUBDOM}\le0.70$ while their
MAC sits in the $0.3$--$0.744$ ambiguous band. One ARTIFACT candidate
($\Omega=0.424067$) shows the same non-selected-block bleed-through
already documented for ffp1's \texttt{art\_d}: its other block
independently reads $\mathrm{MAC}=0.7699$, $\mathrm{SUBDOM}=0.8343$,
both above their respective bars --- confirmation, not a new finding,
that max-over-blocks is not a safe alternative rule at this geometry
either.

\section{Cluster job-number provenance}\label{sec:provenance}
The main text states every result in prose, with no job number in its own
running text; this section is the audit trail for the handful of findings
the main text also flags with a footnote marker but does not itself
enumerate by job. Each entry below is keyed to the main-text location
that cites this section.

\begin{itemize}\itemsep2pt
\item \textbf{\S4, split/anchor/adaptive enhancement.} Job 2409241
recovered 24 of 26 tabulated modes, including the closest pair discussed
in the main text; jobs 2410035/2410044 independently confirmed the two
remaining misses as genuine, distinct roots. The required second-geometry
blind pass (in-plane orthotropic), job 2410036, recovered only 1 of 8
known targets; a follow-up that escalated the search basis, job 2428927,
left the same mixed pattern, ruling out a basis-size explanation.
\item \textbf{\S5, extensional residual portability.} Job 2421896
confirmed the flexural screen unchanged at $\mathrm{dps}=25/40/60$ on the
primary and two further geometries. Job 2421904 showed the extensional
residual ratios still overlap at $r_0/2b=2.0$, the basis for the
one-sided-only claim.
\item \textbf{\S5, FF-P1 own-root separation.} Jobs 2409212/2409213
evaluated the complete FF-P1 REAL-like/ARTIFACT-like set at each
candidate's own converged root at $n=20$ and $n=28$.
\item \textbf{\S6.3, the 1114.9/1115.3~Hz pair.} Job 2428351 isolated the
two golden-section minima in brackets that cannot reach each other; job
2434752 ran MAC against the two isolates.
\item \textbf{\S6.3, candidate 2.828013's basis-sensitive self-parity.}
Jobs 2408715, 2409032, 2409058.
\item \textbf{\S6.4, geometry-sweep recapture.} Job 2416746 is the full,
current-\texttt{SOLVER\_VERSION} recapture of all four sweep geometries
behind Table~\ref{tab:geomsweep}; jobs 2428355--2428358 are the flexural
residual-screen pass over its 24 out-of-plane keys.
\item \textbf{\S6.5, FE mode 9 / FE\#10 neighbor.} Job 2428350 re-derived
the fidelity of the FE\#10 neighbor (reproducing job 2410396's depth);
job 2410396 is the source per-candidate adjudication data for FE\#23; job
2421895 is the unconstrained scan that found no second minimum for
FE\#23.
\item \textbf{\S6.5, crossed-geometry table.} Jobs 2421898/2421899
(solver capture), 2421900/2421901 (ANSYS extracts), 2422897 (merge),
2422898 (adjudication).
\item \textbf{\S6.6, MAC-bar calibration robustness.} Jobs 2431872,
2431873 and 2432008 produced the three-geometry calibration table; job
2432028 independently reproduced job 2432008 bit-for-bit against a
transcribed gate table; job 2431917 re-evaluated the one off-root
candidate at its own root instead of its catalogued seed.
\item \textbf{\S6.6, C-class 80-mode extract.} Jobs 2445267--2445274
rescored all eight $2\Theta/\pi=1.25$ and $1.50$ keys against an
80-elastic-mode PLANE183 extract (modes $4$--$83$). Every key printed
\texttt{C80-C\_SURVIVES}: 0 tight MAC-REAL at threshold $0.744$
(reconstructed $82$--$115$/$82$--$115$ per key). The tabulated 24-key
sweep remains the 57-mode Rank~14 extract; this pass only tests whether
a deeper list creates a REAL side.
\item \textbf{\S6.7, SUBDOM cluster confirmation.} Job 2445746 scored the
curated candidate lists at the three non-primary geometries (all
A\_CONFIRMS, zero \texttt{B\_OVERLAP}); job 2444987 is the primary
geometry's own curated run (15/15, A\_CONFIRMS) and the source of the two
\texttt{PROBE} points that prospectively predicted the pairing-defect
resolution below. Jobs 2445957/2445958 are the full blind-adjudication-set
test, split across independent SLURM array tasks (8-way at
$r_0/2b=2.0$, $2\Theta/\pi=1.0$, 75 candidates; 5-way at
$r_0/2b=2.0$, $2\Theta/\pi=0.5$, 47 candidates; 122 total, the
count cited in main-text \S6.7) and combined by a dedicated merge
script that recomputes each row's verdict from its own stored gate
record. The ANSYS extraction of FE
mode 35's eigenvector at 682.741~Hz (mesh64/mesh96, cross-converged MAC
$0.9991$) is unnumbered in this log (a cluster ANSYS job, not a solver
probe). Job 2446007 is the first MAC cross-check attempt, later found to
have a block-selection bug (both candidates silently reconstructed the
same parity block); job 2446128 is the corrected re-run, with an in-run
fidelity gate proving the fix reproduces the original unforced behavior
bit-for-bit, giving the cited $\mathrm{MAC}=0.9944$/$0.0000$ result.
Job 2446328 forced both blocks at the three Rank-8 leftovers (FE modes
32/50/53) against existing fullblock eigenvectors: FE 53 is a pairing
artifact (both MAC below $0.3$); FE 32's other block matches at
$0.8528$; FE 50 is ambiguous (best nearby MAC $0.4963$, job 2446900).
A later $\Omega$ scan resolved FE 53 as a third pairing defect with a
replacement root: job 2448570 found a SUBDOM-real local $\sigma_{\min}$
minimum on the matching parity block at $\Omega=2.265637$; job 2448862
scored $\mathrm{MAC}=0.9971$ against FE~53 there; job 2448978 found the
standard unconstrained golden-section primitive unsafe on this window
(a non-matching-block structure dominates nearby); job 2449358 therefore
ran a block-constrained golden section on the matching block's own
$\sigma_{\min}$, converging to $\Omega^{*}=2.265896$; job 2449466
re-scored $\mathrm{MAC}=0.9978$ against FE~53 at that polished point
and $\mathrm{MAC}=0.0000$ against the non-matching block.
Job 2446899 scored SUBDOM on both blocks at FE 32's $\Omega$: the
matching block reads $0.7136$. Job 2446521 is the r200/a50 FE mode 16
MAC (both blocks below $0.3$, pairing artifact). Job 2446902 is FE
modes 26/30 at that geometry (FE 26 other-block $\mathrm{MAC}=0.9302$;
FE 30 both below $0.3$). Job 2446898 recomputed four labelled FF-P1
points at $\mathrm{dps}=60$; SUBDOM agrees with the $\mathrm{dps}=40$
table to the printed digits. Job 2446965 is a duplicate of 2446902.
Jobs 2447415--2447418 recomputed the full labelled set at all four
geometries at $\mathrm{dps}=60$ (unique checkpoints); SUBDOM matches
the $\mathrm{dps}=40$ table to the printed digits. Job 2447410 found no
FE owner for leftover F6/FE mode 50 in modes $4$--$60$ (best MAC
$0.4963$); job 2447861 scored both blocks SUBDOM~$=0.0894$/$0.4813$.
Job 2448321 scanned $\Omega$ across this window and found a second
structure near $2.144515$; job 2448496 scored $\mathrm{MAC}=0.9974$
against FE\#50 there; a golden-section refinement (matching
\texttt{find\_modes\_sigmin}'s own Phase 2 primitive) polished it to
$\Omega=2.144291$ (SUBDOM~$=0.9974$, $\mathrm{MAC}=0.9993$ on recheck),
resolving the pairing.
A stored-data-only search (jobs 2446899, 2446902, 2447411, 2447412,
2445746; no new solves) found no block-pick rule, other than
\texttt{RULE\_LOCALMIN\_RESID} itself or max-over-blocks, that
reaches FE mode 32 or 26 without also certifying a known artifact
(ffp1 candidate \texttt{art\_d}) REAL; both remained open misses at
$n_{\mathrm{dofs}}=20$. Jobs 2448222/2448223 then tested basis
enlargement: FE mode 32 stays a clean miss at $n_{\mathrm{dofs}}=24$,
$x_{\max}=40$ and $28/60$ (its relocated well is an unrelated mode,
$\mathrm{MAC}\approx0.06$); FE mode 26 shows a basis-fragile signal
at $n_{\mathrm{dofs}}=24$ (\texttt{RELOCATED\_REAL},
$\Omega=2.076435$) that does not reproduce at $28/60$ within a
$\pm0.04$ window. Job 2448320 widened the $n_{\mathrm{dofs}}=28$,
$x_{\max}=60$ window to $\pm0.10$ and found a deeper well at
$\Omega=2.112435$ plus a second, shallower minimum nearby, leaving
the shape (one valley or two) unresolved. Job 2448501 ran a finer
$0.002$-step connecting scan across both and resolved it: a single
continuous $\sigma_{\min}$ valley with its true minimum at
$\Omega=2.110435$ (SUBDOM~$=0.9985$, $\mathrm{MAC}=0.8919$, a
single clean qualifying local minimum). A golden-section refinement
(\texttt{golden\_refine\_ip\_r200a50\_m26\_v1.py}, same primitive
and bracket-construction convention as the FE-50 refinement above)
polished this to $\Omega=2.110860724725\ldots$, gated against an
already-validated same-basis catalogue root (r200/a50 $s_1$, drift
$1.1\times10^{-3}$ against a $5\times10^{-3}$ bar) and against its
own scan-grid estimate (shift $4.3\times10^{-4}$ against a $2\times
10^{-3}$ bracket half-width). The MAC/SUBDOM recheck at this polished
value first read $\mathrm{MAC}=0.4578$, a sharp, suspicious drop from
the scan-grid's $0.8919$ over a $4\times10^{-4}$ shift in $\Omega$;
tracing it found that the recheck script's load of
\texttt{probe\_ip\_pairing\_defect\_r200a100\_v2.py} had not set the
\texttt{MAC\_N\_DOFS}/\texttt{MAC\_XMAX} environment variables that
module reads at import time to bind its own default basis, so the MAC
reconstruction had silently fallen back to $n_{\mathrm{dofs}}=20$,
$x_{\max}=20$ instead of the enlarged basis in use --- the same class
of import-time default-binding pitfall documented earlier in this
section for \texttt{XMAX0}. Setting those variables before the module
load, and confirming the fix reproduces job 2448501's stored
$\mathrm{MAC}=0.8919$ bit-for-bit at the identical scan-grid $\Omega$
before trusting the polished-$\Omega$ number, gives the corrected
recheck: SUBDOM~$=0.9991$, $\mathrm{MAC}=0.8906$, a single clean
qualifying local minimum, both comfortably above their respective
bars -- the cited FE mode 26 recovery (main text \S6.7).
Job 2447411 scored FE 26/30 SUBDOM on both blocks (q1 $0.8728$ at FE 26;
q0 $0.7942$ at FE 30). Jobs 2447412--2447414 are the labelled
SUBDOM-vs-MAC bake-off (REAL/HELD all same-block; $t6$/$t8$/$s8$
ARTIFACT-ambiguous). Jobs 2447419/2447420/2447863 are the KILL2
$n=28$/$x_{\max}=60$ scan: catalogue $\Omega$ remains a local minimum
at SUBDOM~$=0.1857$; the neighbouring genuine root $1.987180$ relocates
to $1.981966$ (SUBDOM~$=0.9886$). A stored-checkpoint re-score at
$\gamma=4$ is $45/49$ and $25/28$ with zero duplicates (jobs
2445957/2445958 rows). Jobs 2447871--2447874 are the labelled set at
$n_{\mathrm{dofs}}=24$, $x_{\max}=40$: local-minimum points still
separate at all four geometries (the in-container r200/a100 window is
reproduced); r200/a50's printed \texttt{B\_OVERLAP} is a non-root
negative control sitting on a well that exists at this basis, not a
labelled-root miss. Jobs 2447875--2447878 are the labelled set at
$n_{\mathrm{dofs}}=28$, $x_{\max}=60$: KILL2 remains a local minimum at
SUBDOM~$=0.1857$; the other three geometries confirm on the points that
are still roots. Jobs 2447879--2447882 repeat $n=28$ at each geometry's
own production $x_{\max}$ (KILL2 is REAL at $x_{\max}=20$; the miss is
the $x_{\max}=60$ setting). Jobs 2447883/2447884 are the full blind
adjudication set at $n=28$/$x_{\max}=60$, merged from the same array
pattern as 2445957/2445958: $28/31$ and $14/15$ FE modes that still have
a local-minimum claimant, zero duplicates accepted (many catalogue
$\Omega$ are \texttt{C\_MOVED} and excluded). Job 2447885 extracted
r200/a50 fullblock eigenvectors for modes $4$--$60$ (57 files, mesh96);
job 2447886 found FE mode 30's non-picked block matches FE mode 28 at
$\mathrm{MAC}=0.7709$, not FE mode 30. Job 2433080 is the 80-mode
PLANE183 extract (modes $4$--$60$, mesh96) supplying the fourth
geometry's own SUBDOM-vs-MAC bake-off, above; that bake-off itself ran
in the diagnostic container, not as a separate SLURM job.
\item \textbf{\S6.8, Wang orthotropic table.} Job 2408719 (targets 5--8),
resuming targets 1--4 from job 2408466; an in-run fidelity gate on the
isotropic reduction was confirmed before either search ran.
\item \textbf{\S6.9, Qin table.} Job 2409260 is the original $n=12$
search. Jobs 2421062/2421063 are the Qin-targeted $n=20$/$28$ re-search
with an enlarged $x_{\max}$, each realizing 18 degrees of freedom.
Job 2421905 isolated the second nearby dip
at $\Omega_{\mathrm{Qin}}=4.4717$. Job 2434753 ran MAC on the two
isolated dips against their respective FE targets.
\item \textbf{\S6.10, Shi et al.\ out-of-plane table.} Job 2414020
produced the four directly refined roots and set the acceptance criteria
(depth $\le-3.5$, error below 2\%, bracket-shift replicas within 0.1\%)
reused for the remaining two. Job 2421061 resolved the 53.649 entry's
$0.59\%$ bracket-shift walk into two stable, genuine roots. Job 2411852
confirmed the 38.428 entry's $n=20$ near-miss was not a scan-resolution
artifact. Job 2421581 refined 38.428 at $n=28$/$36$ to a shift-stable
root.
\end{itemize}

\FloatBarrier
\begin{thebibliography}{9}\small
\bibitem{seok1} J. Seok, H.F. Tiersten, ``Free vibrations of annular sector
cantilever plates. Part 1: out-of-plane motion,'' J. Sound Vib. 271 (2004)
757--772.
\bibitem{seok2} J. Seok, H.F. Tiersten, ``Free vibrations of annular sector
cantilever plates. Part 2: in-plane motion,'' J. Sound Vib. 271 (2004)
773--787.
\bibitem{allemang2003} R.J. Allemang, ``The modal assurance criterion --
twenty years of use and abuse,'' Sound Vib. 37(8) (2003) 14--21.
\end{thebibliography}
\end{document}
